\chapter{Introducción}
\label{cap:introduccion}




\section{Motivación}
\label{sec:motivacion}

La mayoría de estudiantes enfrentan dificultades para gestionar adecuadamente su hogar debido a la falta de tiempo, ya que gran parte de su jornada
está dedicada a sus estudios. Esta situación convierte la organización y la realización de tareas domésticas en un problema, pues estas actividades requieren 
una planificación cuidadosa y tiempo para llevarlas a cabo. Por otro lado, la gestión de los gastos compartidos puede resultar especialmente compleja en un entorno 
de estudiantes, ya que es común que cada miembro del hogar se encargue de diferentes pagos. Esto puede generar confusiones y equivocaciones, 
respectoa a quién ha pagado ciertos gastos y qué cantidades están pedientes, lo que agrava aún más la dificultad de una correcta
administación del hogar.

Después de 5 años viviendo de manera independiente, he experimentado de primera mano las dificultades de gestionar un hogar compartido con otros estudiantes o trabajadores. En numerosas ocasiones 
el tiempo libre del que disponíamos no se ha destinado a la organización del hogar. Aunque las tareas se asignaban semanalmente, no siempre se cumplían, ya sea
por falta de tiempo y o por descuidos. Además, la administación de los gastos ha sido, en muchas ocasiones, caótica, ya que no siempre se llevaba un control adecuado sobre quién ha pagado qué,
y algunas personas solían olvidar devolver el dinero a tiempo.  

Actualmente no existe ninguna aplicación que ofrezca una solución integral tanto para la gestión de las tareas domesticas como para la gestión de los gastos. Las aplicaciones
disponibles abordan estos aspectos por separado. Además la mayoría de ellas están orientadas a hogares familiares donde la distribución de las tareas no se gestiona de la misma forma que 
en los pisos de estudiantes. Por eso es fundamental desarrollar una aplicación que integre todas estas funcionalidades, dirigida específicamente a estudiantes o trabajadores que compartan vivienda. 

\section{Objetivos del proyecto}
\label{sec:objetivos}

El objetivo general de este proyecto es facilitar a los estudiantes la organización y gestión eficiente de su hogar, ofreciendo soluciones que aborden las dificultades 
que puedan surgir en la planificación y distribución de responsabilidades. Para poder alcanzar este objetivo, podemos diferenciar una serie de objetivos específicos:

\begin{itemize}
	
	\item \textbf{Objetivo 1: Estudio de aplicaciones similares}
	\begin{itemize}
		\item \textbf{Descripción:} Estudio del mercado actual de aplicaciones similares con el fin de identificar funcionalidades innovadoras que puedan ser incorporadas en nuestra aplicación.
	\end{itemize}

	\item \textbf{Objetivo 2: Análisis de las tecnologías a utilizar}
	\begin{itemize}
		\item \textbf{Descripción:} Investigar cuáles son las tecnologías más adecuadas para la implementación del proyecto.
	\end{itemize}
	
	\item \textbf{Objetivo 3: Desarrollo e implementación de la aplicación}
	\begin{itemize}
		\item \textbf{Descripción:} Analizar, desarrollar y realizar pruebas en la aplicación con el fin de obtener los mejores resultados.
	\end{itemize}
	
	\item \textbf{Objetivo 4: Validación de la aplicación por usuarios reales}
	\begin{itemize}
		\item \textbf{Descripción:} Evaluar la aplicación en un entorno real para verificar si satisface las expectativas del usuario.
	\end{itemize}
    \item \textbf{Objetivo 5: Análisis de usabilidad}
    \begin{itemize}
        \item \textbf{Descripción:} Realizar un análisis de usabilidad para identificar posibles mejoras en la interfaz de usuario.
    \end{itemize}
\end{itemize}

\section{Estructura de la memoria}
\label{sec:estructura}

A continuación, se ofrece una visión general de la estructura y organización de la memoria. El contenido se ha dividido en los siguientes apartados:

\begin{enumerate}
	\item \textbf{Resumen.} Un resumen del proyecto para ofrecer una visión general del mismo.
	\item \textbf{Introducción.} En este apartado se presenta la motivación del proyecto y los objetivos que se pretenden alcanzar.
	\item \textbf{Planificación.} Descripción de la planificación del proyecto, que incluye las diferentes fases que se han llevado a cabo durante su desarrollo, así como una estimación de los recursos económicos necesarios.
	\item \textbf{Estado del arte.}. Estudio sobre los principales temas relacionados con este proyecto, cómo analisis de aplicaciones similares, temas legales de protección de datos, etc.
	\item \textbf{Análisis de tecnologías a utilizar.} Análisis de las tecnologías utiizadas para la implementación del proyecto. Para cada una de dichas tecnologías se describen sus características y se justifica su elección.
	\item \textbf{Análisis de la aplicación.} Se detalla el análisis de la aplicación, incluyendo los requisitos funcionales y no funcionales, así como los casos de uso.
	\item \textbf{Diseño.} Descripción del diseño de la aplicación, incluyendo la arquitectura, el diseño de la base de datos y la interfaz de usuario.
	\item \textbf{Desarrollo e implementación de la aplicación.} Se detalla el proceso de desarrollo de la aplicación. Se realizará siguiendo una metodología ágil y se detallarán sus iteraciones.
	\item \textbf{Validación de la aplicación por usuarios reales.} Evaluación de la aplicación en un entorno real para verificar si satisface las expectativas del usuario.
	\item \textbf{Conclusiones y trabajo futuro.} Se presentan las conclusiones obtenidas a lo largo del proyecto y se proponen posibles líneas de trabajo futuro.
	\item \textbf{Bibliografía.} Referencias bibliográficas utilizadas en la elaboración de la memoria.

\end{enumerate}


\newpage